% STATUS: verified
\chapter{The system and the course}\label{ch:course}

The \product{} was built as the project of a course, and the course's method is its
specification. This chapter sets the two side by side: what the course's three assignments ask
for and where each answer is in the product; which of the course's concepts became which part
of the system; where the system follows the text of a source rather than a summary of it; and
what the course teaches that the system does not do.

\section{The course and its sample organisation}\label{sec:course:course}

The course is \lt{Finansų teisė ir finansų įstaigų veiklos atitikties valdymas} --- financial
law and the compliance management of financial institutions --- in the FinTech master's
programme of Vilnius University Kaunas Faculty, autumn 2026. Its lectures run from the
structure of the financial markets and the licensing of each kind of institution, through the
place of the compliance function in internal governance and the building of a compliance
management system, to the requirements on money-laundering prevention and on digital
operational resilience, and to liability for breaches. Its independent work is done on one
sample organisation, \demo{}: a simulated crypto-asset service provider of about forty people,
published under a Creative Commons licence with 26 policies --- thirteen of an information
security management system and thirteen on its financial-technology business
(\cref{ch:company}).

The course also describes a role: the technologist who builds tools for legal and compliance
work. The duties it lists for that role read like a checklist for a system of this kind ---
prompts whose answers are bound to a schema, answers grounded in retrieved sources, regression
suites that include a model judging a model, kept sets of ground-truth cases, and audit
trails. \Cref{sec:course:duties} goes through them.

\section{The three assignments}\label{sec:course:assignments}

Each assignment asks a student to obtain a set of documents about \demo{} with the help of a
language model. The system treats the second and the third as products: something generated
from the company's own documents, checked by code against the form the assignment prescribes,
and left for a person to edit and approve; the first it answers in the chat and files on the
company. \Cref{tab:course:assignments} gives the mapping; the
paragraphs after it say what the form is and how it is enforced.

\begin{table}[H]
\caption{The course's assignments and where each is in the product}\label{tab:course:assignments}
\begin{tabularx}{\textwidth}{@{}L{22mm}YY@{}}
\toprule
\hd{Assignment} & \hd{What it asks for} & \hd{Where it is in the product} \\
\midrule
1, first part & Classify an organisation: the financial market it works in, the institutes of
  financial law that govern it, and the licence or list entry it needs. &
  A conversation with the chat assistant, which answers from the corpus with citations and
  files the result on the company's profile (\code{financial_market}) through its tools
  (\cref{ch:chat,ch:company}). \\
1, second part & A role profile --- for example the compliance officer's --- drawn from the
  company's own documents. &
  A role of the company (\code{company_roles}); a chat thread can adopt it, and its text then
  enters the chat's prompt as the variable \code{ROLE_CONTEXT} (\cref{ch:chat}). \\
2 & A compliance policy in eight parts and the lists of missing second-level policies and
  third-level procedures; then the positions found in the documents, an organisational
  structure, the compliance process in steps, and a responsibility matrix. &
  The work products \code{policy} and \code{governance}, exported as PDF and, for the policy,
  Word (\cref{ch:programme}). \\
3 & Criteria of likelihood and impact, a five-by-five matrix with bands and responses, at
  least five risk scenarios in a set sentence, and the legal risk register. &
  The work product \code{risks}, exported as PDF and Excel (\cref{ch:programme}). \\
\bottomrule
\end{tabularx}
\end{table}

\paragraph{Assignment 2: the policy.} The course teaches a compliance policy in eight parts,
and the product's policy has exactly those eight, in that order: the introduction with
purpose, scope and fit; management's commitment and its link to the company's values; the
mandate of the compliance function; the commitment to meet requirements and to improve
continually; roles and responsibilities; the core compliance processes; the reporting of
breaches; and the consequences of non-compliance. Around the policy stands the course's
hierarchy of documents --- the policy itself on the first level, thematic policies on the
second, procedures on the third --- and the product places every document of the company on a
level and lists the second- and third-level documents that are missing, each with the legal
basis that calls for it where a provision does. Code refuses to let a policy be approved while
one of the eight parts has fewer than 150 characters or an article it names does not exist in
its act.

\paragraph{Assignment 2: structure, process and the responsibility matrix.} The assignment's
second half is a chain: the positions the documents name, the roles that are missing, a
structure, a process, a matrix. The product follows the chain in one work product and enforces
the assignment's own rule on the matrix by code: \emph{every step has exactly one accountable
role and at least one responsible role}. It also holds the structure to the three lines of
defence the course teaches --- the second line may not sit under the first, and internal audit
reports to a governing body alone --- requires the process to have at least five steps, and
checks the clauses of ISO~37301 and the paragraphs of the Bank of Lithuania's guidelines on the
compliance function that the steps say they rest on.

\paragraph{Assignment 3: the risk register.} The course prescribes the criteria closely, and
the product's constants are the course's: likelihood on five levels over four dimensions (the
state of internal control, the clarity of the law and the attention of the supervisor, the
counterparties, the frequency), impact on five levels over four (financial, supervisory and
licence, reputation, reach), the score as their product, and five bands --- 1 to 3
\lt{nereikšminga}, 4 to 6 \lt{toleruotina}, 8 to 12 \lt{vidutinė}, 15 and 16 \lt{didelė},
20 to 25 \lt{kritinė} --- each with the responses it allows. A scenario is written in the
course's sentence, \enquote{\lt{Tikėtinas} [event], \lt{kylantis iš} [source], \lt{gali
sukelti} [consequences], \lt{kurios paveiks} [objective]}, and typed by the six kinds of legal
risk of ISO~31022 Annex~A; the product requires at least five scenarios over at least four of
the six kinds, each tied to a unit and a policy of the company. The numbers are where the
product departs from simply asking a model: the model gives the eight levels and its reasons,
and code computes the likelihood, the impact, the score and the band (\cref{ch:programme}).
The register exports in the course's nine columns.

\begin{humanbox}{What the product does not do for the student}
It does not approve its own work. Each of the three work products comes out as a draft with
the rules it meets and the rules it breaks listed beside it; a person edits it and a person
approves it, and the PDF's cover says whether anyone has. An assignment handed in from the
product is therefore a reviewed draft or an unreviewed one, and says which.
\end{humanbox}

\section{The course's concepts in the system}\label{sec:course:concepts}

Beyond the assignments, the course's lectures on monitoring and on the technologist's role
supplied most of the system's shape. \Cref{tab:course:concepts} lists the concepts and the
parts that embody them.

\begin{xltabular}{\textwidth}{@{}L{40mm}YL{20mm}@{}}
\caption{Concepts the course teaches and the part of the system that embodies each}\label{tab:course:concepts}\\
\toprule
\hd{The course teaches} & \hd{The system does} & \hd{Chapter} \\
\midrule
\endfirsthead
\toprule
\hd{The course teaches} & \hd{The system does} & \hd{Chapter} \\
\midrule
\endhead
\bottomrule
\endlastfoot
A finding has four parts: the condition, the criterion, the cause and the consequence. &
  A finding of a report holds all four. The condition is the assessment and its status; the
  criterion is the obligation with the provision behind it; the consequence is joined by code
  from the act's own sanction provisions; the cause is a field only a person fills, because
  documents show that something is missing and never why. & \cref{ch:report} \\
A report answers three questions: what, so what, now what. &
  The findings say what; the consequence and the severity of the obligation say so what; the
  action plan, which covers every open finding, says now what. The PDF is ordered from the
  conclusion down. & \cref{ch:report} \\
A person approves the scope of the work, and the outcome of each task. &
  A report stops after deciding which acts apply and waits for a person to approve the list
  before the costly part (a person may choose to let it run through, and the report then
  records that nobody approved its scope); a person reviews findings and signs the finished
  report off. &
  \cref{ch:report} \\
Findings and actions are kept in a register with the states open, in progress, late and
closed, and an action is closed only after a re-test. &
  The action register. An owner may call an action done, but only a later signed-off report
  --- the re-test --- that finds its obligations met closes it, unless a person closes it by
  hand with a note that says why. Late is derived from the due date, never set by hand. &
  \cref{ch:register} \\
The company's own rules are criteria too, not only the law. &
  A report can assess whether the company's procedures carry out what its own policies
  require, as a group of findings counted apart from the legal ones. & \cref{ch:report} \\
The three lines of defence. &
  The governance work product places every body, function and unit on a line and refuses to
  approve a structure that puts the second line under the first. & \cref{ch:programme} \\
Risk is likelihood times impact, on stated criteria. &
  The risk work product, with the score and the band computed by code. & \cref{ch:programme} \\
\end{xltabular}

\section{The technologist's five duties}\label{sec:course:duties}

The course's first lecture describes the technologist of legal operations for AI and
compliance systems as the person who builds, tests and deploys prompt-based validation
processes --- processes that audit a company's documents against its internal policies and
against the law --- and it names five duties. They are the most direct statement of what the
course expects of a system like this one. \Cref{tab:course:duties} holds the system against
each duty, including where it stops short.

\begin{xltabular}{\textwidth}{@{}L{34mm}YY@{}}
\caption{The lecture's five duties of the technologist, and the system}\label{tab:course:duties}\\
\toprule
\hd{The duty, as the lecture puts it} & \hd{What the system does} & \hd{Where it stops short} \\
\midrule
\endfirsthead
\toprule
\hd{The duty, as the lecture puts it} & \hd{What the system does} & \hd{Where it stops short} \\
\midrule
\endhead
\bottomrule
\endlastfoot
\textbf{Prompt engineering and schema enforcement:} precise, structured prompts, constrained
  by JSON schemas, that check company documents and flag non-compliant clauses or factual
  inconsistencies. &
  Twenty-one of the 23 prompts answer in a strict JSON schema, which the provider is asked to
  enforce and the assistant service checks again (\cref{ch:models,ch:prompts}). The gap
  analysis gives every obligation a status and quotes the evidence it found word for word; a
  document's inventory lists
  its weaknesses and its references to the law; the citation check flags a reference to an
  article that does not say what the document claims (\cref{ch:report,ch:company}). &
  The chat's two prompts answer in free text. \\
\textbf{Grounding and context verification:} retrieval-augmented prompts that compare
  documents directly with internal policy libraries, standard operating procedures and the
  relevant regulations. &
  The law is taken from the corpus, not from a model's memory, and a provision a model names
  from memory --- a work product's legal basis --- is looked up there: the chat retrieves
  provisions by hybrid search; a report hands the model the provisions themselves, or the
  obligations read from them (\cref{ch:corpus}). The company's whole library goes with every call of the gap analysis ---
  retrieved passages were used first and replaced, because they missed evidence the documents
  held (\cref{ch:report}). A company's own policies can be the criteria too. &
  Retrieval remains the fallback for a library too large to send whole, and the chat still
  works from what it retrieves. \\
\textbf{Automated evaluation and hallucination defence:} regression suites --- a model as
  judge, natural-language inference, unit tests --- that keep measuring the reliability of
  prompts, the accuracy of citations and the handling of edge cases. &
  Six suites of cases, labelled on 12~November 2026 by an AI model and confirmed by no person
  yet, run against the real model through the product's own functions and are held to a
  committed baseline; a second model call judges in three places; some 6\,100 tests run without
  calling a model (\cref{ch:control,ch:testing}). The defence that does not depend on any
  model: every quotation a finding or a work product rests on is located in its document,
  every provision a structured answer names is looked up, every score is computed. &
  No natural-language-inference model is used. The suites measure six of the 23 prompts,
  and are run by hand before and after a change of prompt or model, not continuously. \\
\textbf{Collaboration with experts and golden datasets:} turning complex regulatory
  requirements into deterministic prompt instructions together with compliance officers and
  lawyers, and keeping verified ground-truth datasets. &
  The register of obligations is the requirements in a deterministic form: each duty written
  once, tied to its provision, judged and reviewed, and reused by every report
  (\cref{ch:obligations}). Every status and verdict a prompt may give is defined in the prompt
  (\cref{ch:prompts}). The labelled cases are kept in the repository, in five of the six
  suites with the reason for each label. &
  The labels were made on 12~November 2026 by the AI coding assistant used to build the
  system (Claude), reading the sources; the register that ships is machine-reviewed, judged
  by the judge and not read by a person entry by entry (\cref{ch:obligations,ch:control}). No
  compliance officer or lawyer
  has yet confirmed either. Both are built for a person to overrule any label or entry. \\
\textbf{Auditability and integration:} validation built into the existing systems of legal
  operations --- document management, contract management, process engines --- with full
  event logging and audit trails. &
  Every change by anyone is in the audit log; every model call that answers is logged with
  its prompt version, tokens and cost (the embedding of a search query excepted), and every
  chat turn with its transcript; every report records
  the versions of the prompts and acts it was made from (\cref{ch:security,ch:models}). The
  whole system is reachable through a JSON API (\cref{ch:api}). &
  No connector to a document-management or contract-management system is built: documents
  are uploaded. \\
\end{xltabular}

\section{Where the system follows the text of a source}\label{sec:course:sources}

In two places the system's behaviour comes from the wording of a source itself, and a reader
comparing it with a summary of that source may notice.

\begin{itemize}
  \item \textbf{The Bank of Lithuania's guidelines on the compliance function} are worded as
    recommendations (\lt{turėtų}, \enquote{should}) and are addressed to the kinds of
    institution their third paragraph lists, which do not include crypto-asset service
    providers. The catalogue scopes them accordingly. For a company such as \demo{} a report
    therefore does not assess them on its own; a person may take them into the scope of a
    report on the scope screen, and they are then assessed as good practice and labelled so
    (\cref{ch:report}).
  \item \textbf{Clause numbers of ISO~37301.} The steps of the compliance process cite the
    standard's clauses as the standard numbers them: raising concerns is clause~8.3 and
    investigation processes clause~8.4. A step that cites no clause the standard has, and no
    paragraph of the guidelines, is flagged by the product's own check, as a warning
    (\cref{ch:programme}).
\end{itemize}

\section{What the system adds to the course's method}\label{sec:course:adds}

Four things in the system are not asked for by any assignment and follow from taking the
method seriously as software.

\begin{description}
  \item[The law as a maintained corpus.] The assignments let a model speak about the law. The
    system fetches the acts from their official sources, keeps them in versions, and records
    what each new version changed --- so that it can say which findings rest on a provision
    that has changed since (\cref{ch:corpus,ch:register}).
  \item[A register of obligations.] Between the law and a report stands a reviewed list of
    what each act obliges, extracted once and reused by every report. It is what makes two
    reports of two companies comparable, and what gives a person one place to correct the
    reading of the law (\cref{ch:obligations}).
  \item[The citation audit.] The sample organisation's overview document cites some articles
    of the regulation on markets in crypto-assets by the wrong number. The system checks every
    article a document cites against the corpus and against what the document says the
    article covers (\cref{ch:company}).
  \item[Cost as a design constraint.] Every call the system records is priced; a report shows its estimate before
    it starts; the costly part waits for a person, unless the run was started without the
    scope review; one-time work on the law is separated from
    per-company work (\cref{ch:models,ch:control}).
\end{description}

\section{What the course teaches that the system does not do}\label{sec:course:not}

\begin{itemize}
  \item \textbf{The annual compliance report and a monitoring programme.} The course treats
    both as products of a compliance function. The system has the material for them --- the
    reports, the action register, the history of each obligation --- but does not draft
    either.
  \item \textbf{A maturity assessment} of the compliance management system is not a product of
    the system.
  \item \textbf{The ISO standards themselves.} The system can read a standard into clauses and
    extract its requirements (\cref{ch:corpus,ch:obligations}), and a company can adopt one
    and be assessed against it. The course's copies of the ISO standards are scanned images
    under a teaching licence: the system does not read scans, and the licence does not cover
    loading them. No standard is in the deployed corpus.
  \item \textbf{Operating effectiveness.} The course distinguishes the design of a control
    from its operation. A report judges design, from documents; it marks where operation
    should be verified and leaves the verification to a person.
  \item \textbf{The cause of a finding.} Deliberately left to the reviewer.
\end{itemize}

The limitations of the AI paths themselves --- what is not measured, what a judge cannot
establish, how the shipped register of obligations was reviewed --- are gathered in
\cref{ch:control}.
